Os cogumelos constituem uma relevante fonte natural de nutrientes, destacando-se pelo elevado teor de proteínas, vitaminas e compostos bioativos benéficos à saúde humana. No mercado global de cogumelos comestíveis, as doenças fúngicas representam uma das principais causas de perdas produtivas. Esses organismos são cultivados predominantemente em ambientes fechados, sob condições ambientais controladas, o que favorece a adoção de programas integrados de manejo de doenças. Tais programas combinam o uso de fungicidas químicos, alternativas de biocontrole e práticas agronômicas adequadas, permitindo reduzir a incidência de patógenos e garantir o manejo sanitário das culturas. Dessa forma, diversos métodos de controle são continuamente aplicados com o objetivo de minimizar perdas e assegurar a qualidade do produto final.\cite{Bellettini2017}

Com base nesse cenário, a aplicação de tecnologias emergentes no setor agrícola possibilita o aumento da eficiência produtiva, a redução de perdas e a melhoria da qualidade dos alimentos e fatores diretamente relacionados à segurança alimentar. Além disso, o monitoramento contínuo de variáveis ambientais, como umidade e temperatura, permite otimizar o uso de recursos naturais, reduzir o desperdício de insumos e promover uma produção mais sustentável, resiliente e alinhada às demandas atuais por sistemas agrícolas inteligentes e ecologicamente responsáveis.

No contexto brasileiro, o Pleurotus ostreatus (shimeji) destaca-se entre as espécies mais amplamente cultivadas, devido ao seu elevado valor nutricional e à crescente demanda no mercado alimentício. Entretanto, o aumento da produção tem sido acompanhado por uma maior ocorrência de contaminações microbiológicas, resultantes da proliferação de microrganismos oportunistas durante as etapas de cultivo e processamento, podendo ocasionar perdas totais nas colheitas. As infestações por pragas e doenças comprometem significativamente a qualidade e a composição do produto final, reduzindo a competitividade do produtor no mercado e impactando negativamente a rentabilidade e a sustentabilidade da produção.\cite{Scheffer2023}

Com isso o Vale do Ribeira, localizado no Estado de São Paulo, destaca-se como uma região promissora para a implementação de sistemas baseados em tecnologias inteligentes. Essa região, foco da presente pesquisa, apresenta um constante crescimento na produção de espécies de cogumelos, especialmente o shimeji, o que reforça o potencial de aplicação de soluções tecnológicas voltadas ao monitoramento e controle de variáveis ambientais, contribuindo para o fortalecimento da produção local e o desenvolvimento sustentável do setor.\cite{Urben2017}

Deste modo o presente projeto tem como objetivo o desenvolvimento de um sistema inteligente de monitoramento do cultivo de cogumelos, voltado à prevenção de infecções e infestações por parasitas. A proposta visa integrar sensoriamento ambiental e inteligência artificial (IA) para garantir o acompanhamento automatizado das condições do composto e do desenvolvimento dos cogumelos. O sistema será responsável por coletar e exibir dados em tempo real sobre umidade, temperatura e porcentagem estimada de infecção, obtidos por meio de sensores instalados no ambiente de cultivo, permitindo que o produtor acompanhe remotamente as condições ideais para o crescimento saudável dos fungos.

Além disso, será implementada uma IA baseada em aprendizagem profunda (deep learning). As imagens capturadas no ambiente serão processadas e utilizadas tanto para o treinamento da IA quanto para o monitoramento contínuo do cultivo. O modelo será treinado para identificar padrões visuais indicativos de infecções e parasitas, como alterações de cor, presença de manchas e irregularidades no composto. Quando o sistema detectar anomalias que indiquem risco de contaminação, a porcentagem de alerta será automaticamente ajustada e o produtor será notificado por meio do aplicativo, permitindo uma resposta rápida e precisa no manejo das culturas. 

Alinhando-se o presente estudo com a Agenda 2030 e os Objetivos de Desenvolvimento Sustentável (ODS) propostos pela Organização das Nações Unidas (ONU), a presente pesquisa relaciona-se diretamente ao ODS nº 2 — Fome Zero e Agricultura Sustentável. Este objetivo visa erradicar a fome e promover sistemas alimentares produtivos, resilientes e ambientalmente responsáveis. Dessa forma, o desenvolvimento de um sistema inteligente de monitoramento e prevenção no cultivo de cogumelos contribui não apenas para o aumento da produtividade e da segurança alimentar, mas também para a redução do desperdício e o uso eficiente dos recursos naturais, reforçando o compromisso com práticas agrícolas sustentáveis e inovadoras.
